\documentclass[11pt]{article}
\usepackage{amsmath,amssymb,amsthm,mathtools}
\usepackage[margin=25mm]{geometry}
\usepackage[hidelinks]{hyperref}
\usepackage{longtable}
\newtheorem{theorem}{Theorem}
\newtheorem{lemma}[theorem]{Lemma}
\newtheorem{proposition}[theorem]{Proposition}
\newcommand{\M}{\mathcal M}
\newcommand{\F}{\mathcal F}
\newcommand{\ii}{\mathrm i}
\newcommand{\dd}{\,\mathrm d}
\setlength{\emergencystretch}{3em}
\title{The original theta input: exact analytic maps, integer tails, and certified Toda seed data}
\author{}
\date{12 September 2026; certificate completed 13 September 2026}
\begin{document}
\maketitle

This continuation supplies the omitted-tail and ball-arithmetic certificate
for the analytic input of Section 9 of \emph{Source--boundary volume
dynamics for the arithmetic control}. Every theta coefficient, Fourier
factor, conjugation, and complex branch used below is specified. The
original seed has total mass $M_1(0)$; the measure is never divided by
that mass. The result computes analytic input of the source determinant
flow. Its exact connection with a finite zero packet is given below by
Mellin division and the multiplier $h(1/2+\ii t)$.

\section{The original theta function and its Mellin transform}

Put $D=-z\partial_z$, let
$\mathcal S_{\pi/4}=\{z\in\mathbb C\setminus\{0\}:|\arg z|<\pi/4\}$,
and use the restriction of the principal logarithm to this sector. Define
\begin{align}
 \vartheta(z)&=\sum_{n\in\mathbb Z}e^{-\pi n^2z^2},\tag{TI.1}\\
 f_0(z)&=D(D-1)\vartheta(z)
 =2\sum_{n=1}^\infty(4\pi^2n^4z^4-6\pi n^2z^2)e^{-\pi n^2z^2}.
 \tag{TI.2}
\end{align}
On a compact subset of the sector, each differentiated summand is
bounded by a polynomial in $n$ times $e^{-a n^2}$ with $a>0$.
Consequently both series define holomorphic functions there, and all
the displayed derivatives are termwise valid.

\begin{lemma}[The exact theta inversion]
On $\mathcal S_{\pi/4}$,
\begin{equation}
 \vartheta(z)=z^{-1}\vartheta(1/z),\qquad
 f_0(z)=z^{-1}f_0(1/z),\qquad
 f_0(\bar z)=\overline{f_0(z)}. \tag{TI.3}
\end{equation}
For each $0\le\alpha<\pi/4$, every integer $j\ge0$, and every real
$A>0$, $D^jf_0(z)=O(|z|^{-A})$ as $|z|\to\infty$ and
$D^jf_0(z)=O(|z|^A)$ as $|z|\to0$, uniformly for $|\arg z|\le\alpha$.
\end{lemma}
\begin{proof}
For $x>0$, the Fourier transform with kernel $e^{-2\pi\ii\xi t}$
of $e^{-\pi x^2t^2}$ is $x^{-1}e^{-\pi\xi^2/x^2}$.
For completeness, the value at $\xi=0$ follows from
$\int_{\mathbb R}e^{-\pi t^2}\dd t=1$, whose square is evaluated in
polar coordinates. Differentiation in $\xi$ and integration by parts
give the differential equation $I'(\xi)=-2\pi\xi I(\xi)/x^2$;
the stated transform is its solution with the stated initial value.
Periodize the Gaussian over integer translates. The resulting smooth
periodic function has the absolutely convergent Fourier series with
coefficients $x^{-1}e^{-\pi n^2/x^2}$. Evaluation at zero gives
$\vartheta(x)=x^{-1}\vartheta(1/x)$. Both sides are holomorphic in
the connected sector, so the identity theorem gives the first equation
in (TI.3).

For the linear involution $(Rf)(z)=z^{-1}f(1/z)$, direct differentiation
gives $DR=R(1-D)$ and hence $D(D-1)R=RD(D-1)$. Applying this identity
to $\vartheta=R\vartheta$ proves the second equation. The real
coefficients in (TI.2) give the third. On a closed subsector,
$\Re(z^2)\ge |z|^2\cos(2\alpha)>0$; the series and each derivative
therefore have Gaussian decay at infinity, with a fixed polynomial
factor. The second equation in (TI.3), differentiated using $DR=R(1-D)$,
transports that bound to zero and proves all the endpoint assertions.
\end{proof}

For functions with these endpoint bounds set
\begin{equation}
 (\M F)(s)=\int_0^\infty F(x)x^{s-1}\dd x,
 \qquad (\F_+u)(t)=\int_{\mathbb R}u(y)e^{\ii ty}\dd y.
 \tag{TI.4}
\end{equation}
The Mellin integral is entire. When $\Re s>1$, the sum of the absolute
integrals of the terms in (TI.2) is finite, because it is bounded by a
constant depending on $s$ times $\sum n^{-\Re s}$. Thus
\begin{align}
 (\M f_0)(s)
 &=\pi^{-s/2}\zeta(s)
 \left(4\Gamma(s/2+2)-6\Gamma(s/2+1)\right)\notag\\
 &=s(s-1)\pi^{-s/2}\Gamma(s/2)\zeta(s)=:g(s). \tag{TI.5}
\end{align}
The first line contains the factor $2$ of (TI.2) and the factor $1/2$
from the substitution $u=\pi n^2x^2$. The second follows from
$\Gamma(z+1)=z\Gamma(z)$, which follows by integration by parts.
Equation (TI.5) gives the entire continuation of its right side. In
particular $g=2\xi$ with the original convention in the source note.
Substituting $x\mapsto1/x$ in (TI.4), using (TI.3), proves
\begin{equation}
 g(1-s)=g(s),\qquad \overline{g(\bar s)}=g(s).
 \tag{TI.6}
\end{equation}

\section{Mellin division, the dagger sign, and the tilted observation}

Let $h(s)=\prod_\rho(s-\rho)^{m_\rho}$ be a monic finite product of
actual zeros of $g$, with the selected orders at most their actual
orders and with the selected multiset invariant under
$\rho\mapsto1-\bar\rho$. The complete-order packets used in the source
note satisfy these conditions. Write $d=\deg h$ and
\begin{equation}
 v_h=g/h,\qquad h^\dagger(s)=\overline{h(1-\bar s)},\qquad
 v_h^\dagger(s)=\overline{v_h(1-\bar s)}.
 \tag{TI.7}
\end{equation}
Removable singularities are filled in. The leading coefficient of
$h^\dagger$ is $(-1)^d$, and its zero multiset is that of $h$; therefore
\begin{equation}
 h^\dagger=(-1)^d h,\qquad v_h^\dagger=(-1)^d v_h.
 \tag{TI.8}
\end{equation}

\begin{proposition}[Construction without an endpoint assumption]
For every $\sigma\in\mathbb R$ the integral
\begin{equation}
 F_h(z)=\frac1{2\pi\ii}\int_{\sigma-\ii\infty}^{\sigma+\ii\infty}
 v_h(s)z^{-s}\dd s \tag{TI.9}
\end{equation}
converges on the sector, is independent of $\sigma$, defines a
holomorphic function there, and has the rapid endpoint bounds in the
first lemma on each closed subsector. It satisfies
\begin{equation}
 \M F_h=v_h,\qquad h(D)F_h=f_0,\qquad
 F_h(1/\bar z)=(-1)^d\bar z\,\overline{F_h(z)}. \tag{TI.10}
\end{equation}
\end{proposition}
\begin{proof}
Rotate the Mellin integral of $f_0$ to the ray $e^{\ii\beta}\mathbb R_{>0}$,
where $|\beta|<\pi/4$. The endpoint estimates make both closing arcs
vanish, and give
\[
 g(s)=e^{\ii\beta s}\int_0^\infty f_0(e^{\ii\beta}x)x^{s-1}\dd x.
\]
For $\sigma$ in any compact interval, the absolute integral on the
right is uniformly bounded. Choosing the sign of $\beta$ to be that
of $t$ proves $|g(\sigma+\ii t)|\le C e^{-|\beta||t|}$.
For sufficiently large $|t|$, the explicit finite product $h$ is
bounded below by $c(1+|t|)^d$ on the same strip. On the remaining
compact rectangle, $v_h$ is bounded because its singularities are
removable. These facts give the same exponential estimate for $v_h$,
with a different constant. It is valid for every
$|\beta|<\pi/4$ and every compact interval of $\sigma$.

In (TI.9), choose $|\beta|>|\arg z|$. Absolute convergence and local
uniform convergence of every differentiated integral follow.
Horizontal contour segments tend to zero by the same estimate, proving
independence of $\sigma$. Multiplication by $s^j$ computes $D^jF_h$;
moving $\sigma$ arbitrarily far to the right or left gives respectively
the asserted infinity or zero bounds.

The change $x=e^y$ turns (TI.9) on the line $\Re s=\sigma$ into Fourier
inversion for $e^{\sigma y}F_h(e^y)$. The exponentially decreasing
function $v_h(\sigma+\ii t)$ and its derivatives have the required
integrability, and inversion gives $\M F_h=v_h$. The Fourier inversion
and Plancherel conventions here can be verified by inserting
$e^{-\varepsilon t^2}$ in the inverse integral: Fubini gives convolution
with the Gaussian approximate identity of integral one; letting
$\varepsilon\downarrow0$ gives inversion, and applying the same
calculation to the inner product gives the factor $1/(2\pi)$ in
Plancherel. Thus no change of Fourier normalization is implicit.
Multiplication by $h(s)$ in (TI.9) is $h(D)$ on the source. Fourier
inversion of (TI.5) proves $h(D)F_h=f_0$.

For the conjugate-linear involution
$({\cal R}F)(x)=x^{-1}\overline{F(1/x)}$, substitution gives
$\M({\cal R}F)(s)=\overline{\M F(1-\bar s)}$. Equations (TI.8) and
Mellin uniqueness imply ${\cal R}F_h=(-1)^dF_h$ on the positive axis.
Both sides extend holomorphically to the sector as
$z^{-1}\overline{F_h(1/\bar z)}$ and $(-1)^dF_h(z)$.
Conjugating the resulting identity gives the last equation in (TI.10).
\end{proof}

Here are the exact Hilbert-space maps. The change of variables
\begin{equation}
 U:L^2(\mathbb R_{>0},\dd x)\longrightarrow L^2(\mathbb R,\dd y),
 \qquad (UF)(y)=e^{y/2}F(e^y) \tag{TI.11}
\end{equation}
is unitary, with inverse $u\mapsto[x\mapsto x^{-1/2}u(\log x)]$.
On the rapid sector functions constructed above define, for real
$|\theta|<\pi/2$,
\begin{equation}
 (E_\theta F)(y)=e^{(y+\ii\theta/2)/2}
 F(e^{y+\ii\theta/2}). \tag{TI.12}
\end{equation}
This map has codomain $L^2(\mathbb R,\dd y)$; it is the analytic
translation of $UF$ on this specified domain. Its norm is not asserted
to equal the untilted norm. Shifting the $y$ contour by $\ii\theta/2$
in (TI.4), with the closing vertical edges vanishing by the proved
endpoint bounds, gives
\begin{equation}
 \F_+(E_\theta F_h)(t)=e^{\theta t/2}v_h(1/2+\ii t). \tag{TI.13}
\end{equation}
Indeed $y=z-\ii\theta/2$ contributes
$e^{\ii ty}=e^{\theta t/2}e^{\ii tz}$. In the source expression
the factor $e^{\ii\theta/4}$ remains present.

Define the original, unscaled density and its mass by
\begin{equation}
 w_h(t)=\frac{|v_h(1/2+\ii t)|^2}{2\pi},\qquad
 M_h(\theta)=\int_{\mathbb R}e^{\theta t}w_h(t)\dd t,
 \qquad \mu_h=M_h(0). \tag{TI.14}
\end{equation}
The exponential estimate above proves that the integral defining
$M_h$ and each of its derivatives converge locally uniformly on
$|\Re\theta|<\pi/2$: for a compact substrip choose
$2|\beta|>\sup|\Re\theta|$. By (TI.11)--(TI.13) and Plancherel,
\begin{equation}
 M_h(\theta)=\int_0^\infty|F_h(e^{\ii\theta/2}x)|^2\dd x
 =2\int_1^\infty|F_h(e^{\ii\theta/2}x)|^2\dd x
 \quad (\theta\text{ real}). \tag{TI.15}
\end{equation}
For the second equality the last identity in (TI.10) gives
\[
 F_h(e^{\ii\theta/2}/x)=(-1)^d e^{-\ii\theta/2}x
 \overline{F_h(e^{\ii\theta/2}x)}.
\]
After taking absolute squares, the factor $x^2$ cancels the Jacobian
of $x\mapsto1/x$. The dagger sign is retained before that operation.
Dagger stability alone gives this endpoint identity. If the selected
multiset is additionally stable under ordinary conjugation, then
$v_h(\bar s)=\overline{v_h(s)}$, so $w_h(-t)=w_h(t)$ and
$M_h(-\theta)=M_h(\theta)$. This proves the precise extra symmetry
needed for evenness.

The typed relation between seed and packet input is also explicit:
\begin{equation}
 F_h\xmapsto{\;h(D)\;}f_0
 \qquad
 v_h(1/2+\ii t)\xmapsto{\;h(1/2+\ii t)\;}
 g(1/2+\ii t),\qquad
 w_1(t)=|h(1/2+\ii t)|^2w_h(t). \tag{TI.16}
\end{equation}
The arrow in the first display means the linear differential map on
the rapid sector-function space; (TI.10) and (TI.13) prove its
commutation with the transform. The scalar equality of densities
holds even at selected line zeros by continuity; it does not require
division of a measured zero value by zero.

\section{The complete double series and an explicit omitted tail}

For $h=1$ we have $F_1=f_0$. Set $c_1=-6$, $c_2=4$ and, for real
$\theta$ first,
\begin{equation}
 C_{mn}(\theta)=\pi(m^2e^{\ii\theta}+n^2e^{-\ii\theta}),\qquad
 \alpha_{rs}=r+s+\tfrac12. \tag{TI.17}
\end{equation}
For $\Re C>0$ define
\begin{equation}
 I_p(C)=\int_1^\infty x^p e^{-Cx^2}\dd x
 =\frac{\Gamma((p+1)/2,C)}{2C^{(p+1)/2}}. \tag{TI.18}
\end{equation}
The power uses the holomorphic logarithm of the right half-plane,
real on its positive axis. To check (TI.18), first substitute $u=Cx^2$
for positive real $C$. Both sides extend holomorphically to the right
half-plane: the integral by dominated convergence and the right side
by the same continued incomplete gamma integral with its ray deformable
inside that half-plane. The identity theorem then gives (TI.18)
there with the stated branch.

\begin{theorem}[The retained factor and the locally uniform tail]
For every $\theta$ with $|\Re\theta|<\pi/2$,
\begin{equation}
 M_1(\theta)=4\sum_{m,n\ge1}\sum_{r,s=1}^2
 c_rc_s(\pi m^2)^r(\pi n^2)^s e^{\ii\theta(r-s)}
 \frac{\Gamma(\alpha_{rs},C_{mn}(\theta))}
 {C_{mn}(\theta)^{\alpha_{rs}}}. \tag{TI.19}
\end{equation}
Let $M_{1,J}$ be exactly this sum with $1\le m,n\le J$, for an integer
$J\ge0$. Fix $0\le\theta_0<\pi/2$ and $B\ge0$, and put
\begin{align}
 d_0&=\pi e^{-B}\cos\theta_0,\qquad L=J+1,
 \qquad \rho_a=e^{-d_0(2a+1)},\tag{TI.20}\\
 H_p(a,d_0)&=\sum_{j=0}^p\binom pj a^{p-j}
 \sum_{b=0}^j \genfrac{\{}{\}}{0pt}{}{j}{b}\,
 \frac{b!\rho_a^b}{(1-\rho_a)^{b+1}},\tag{TI.21}\\
 A_a(d_0,B)&=6\pi e^B H_2(a,d_0)+4\pi^2e^{2B}H_4(a,d_0).
 \tag{TI.22}
\end{align}
Here $\genfrac{\{}{\}}{0pt}{}{j}{b}$ is determined by $\genfrac{\{}{\}}{0pt}{}{0}{0}=1$,
$\genfrac{\{}{\}}{0pt}{}{j}{0}=0$ for $j>0$, $\genfrac{\{}{\}}{0pt}{}{0}{b}=0$ for $b>0$, and
$\genfrac{\{}{\}}{0pt}{}{j+1}{b}=b\genfrac{\{}{\}}{0pt}{}{j}{b}+\genfrac{\{}{\}}{0pt}{}{j}{b-1}$, with out-of-range
indices zero. For every $|\Re\theta|\le\theta_0$ and
$|\Im\theta|\le B$,
\begin{equation}
 |M_1(\theta)-M_{1,J}(\theta)|
 \le \mathcal E_J(\theta_0,B):=
 16A_1(d_0,B)A_L(d_0,B)I_8(d_0(1+L^2)). \tag{TI.23}
\end{equation}
Every quantity on the right is explicit and positive. It tends to zero
as $J\to\infty$. In particular (TI.23) supplies the requested locally
uniform real-tilt bound by taking $B=0$.
\end{theorem}
\begin{proof}
For real $\theta$, substitute (TI.2) into the last expression in
(TI.15). A product contributes $2\cdot2\cdot2=8$ before integration:
the endpoint factor $2$ and the two original theta coefficients $2$.
Equation (TI.18) contributes $1/2$. The remaining factor is precisely
$4$, and the powers of the two dilations give
$e^{\ii\theta r}e^{-\ii\theta s}$. This establishes the coefficients
and phases in (TI.19), subject to the absolute convergence now proved.

For integers $a\ge1$, $p\ge0$, and $x\ge1$,
\begin{align*}
 \sum_{n\ge a}n^p e^{-d_0n^2x^2}
 &\le e^{-d_0a^2x^2}\sum_{j\ge0}(a+j)^p
          e^{-d_0(2a+1)j}\\
 &=e^{-d_0a^2x^2}H_p(a,d_0).
\end{align*}
The inequality uses $(a+j)^2-a^2=2aj+j^2\ge(2a+1)j$ and $x^2\ge1$.
For the equality expand $(a+j)^p$ by the binomial theorem and use
$j^r=\sum_b\genfrac{\{}{\}}{0pt}{}{r}{b}(j)_b$. That last identity follows inductively
from $j(j)_b=(j)_{b+1}+b(j)_b$ and the displayed recurrence.
Differentiating the geometric series $b$ times gives
$\sum_{j\ge0}(j)_b\rho^j=b!\rho^b/(1-\rho)^{b+1}$.
Thus (TI.21) is a finite positive rational function, not an unevaluated
infinite sum.

For complex $\theta$ in the rectangle, let
$f^\pm(x)=f_0(e^{\pm\ii\theta/2}x)$ and let $f_J^\pm$ contain only
$1\le n\le J$ in the original series. Since
\[
 \Re(\pi n^2e^{\pm\ii\theta})
 =\pi n^2e^{\mp\Im\theta}\cos(\Re\theta)\ge d_0n^2,
\]
and $|e^{\pm\ii r\theta}|\le e^{rB}$, the positive bound just proved
implies
\begin{align*}
 |f^\pm(x)|, |f_J^\pm(x)|
 &\le2A_1(d_0,B)x^4e^{-d_0x^2},\\
 |f^\pm(x)-f_J^\pm(x)|
 &\le2A_L(d_0,B)x^4e^{-d_0L^2x^2}.
\end{align*}
Here $x^2\le x^4$ retains a bound for the original $-6$ term;
no coefficient of the exact series is changed. These inequalities
prove absolute integrability of the full expanded double series,
locally uniformly on the rectangle. They also prove
\[
 2\int_1^\infty
 |f^+(x)f^-(x)-f_J^+(x)f_J^-(x)|\dd x
 \le16A_1(d_0,B)A_L(d_0,B)
 \int_1^\infty x^8e^{-d_0(1+L^2)x^2}\dd x,
\]
because the product difference is
$(f^+-f_J^+)f^-+f_J^+(f^--f_J^-)$.
The right side is (TI.23). For fixed $d_0>0$, formula (TI.21) bounds
$H_p(L,d_0)$ by a fixed constant times $(1+L)^p$: its denominators
are bounded away from zero uniformly for $L\ge1$. Moreover
$I_8(b)\le e^{-b}\int_1^\infty x^8e^{-d_0(x^2-1)}\dd x$ for
$b\ge d_0$. Hence (TI.23) tends to zero exponentially in $L^2$.

For complex $\theta$, $C_{mn}(\theta)$ stays in the right half-plane
by the displayed inequality. The integral
$2\int_1^\infty f^+(x)f^-(x)\dd x$ is holomorphic on the whole strip.
It equals $M_1$ on the real interval by (TI.15) and conjugation of
$f_0$. Holomorphic uniqueness proves equality on the strip. The same
locally uniform bounds justify (TI.19) there and finish the proof.
\end{proof}

The tail controls differentiated moments as well. If the closed disc
$\{z:|z-\theta|\le r\}$ lies in the rectangle of the theorem, Cauchy's
integral formula gives, for every integer $j\ge0$,
\begin{equation}
 |M_1^{(j)}(\theta)-M_{1,J}^{(j)}(\theta)|
 \le \frac{j!}{r^j}{\cal E}_J(\theta_0,B). \tag{TI.24}
\end{equation}
The bound includes arbitrary finite derivative orders used by a Hankel
determinant. It does not assert a useful estimate when the derivative
order or determinant dimension itself grows without bound.

\section{The second moment and the first tensor coefficient}

Let $H=(D-1/2)f_0$. With $y=\pi n^2x^2$ and
$P(y)=8y^2-12y$, literal differentiation gives
\begin{align}
 (D-1/2)(P(y)e^{-y})
 &=\left(-2yP'(y)+2yP(y)-P(y)/2\right)e^{-y}\notag\\
 &=(16y^3-60y^2+30y)e^{-y}. \tag{TI.25}
\end{align}
Thus
\begin{equation}
 H(x)=\sum_{n\ge1}\sum_{r=1}^3d_r(\pi n^2)^r x^{2r}e^{-\pi n^2x^2},
 \qquad(d_1,d_2,d_3)=(30,-60,16). \tag{TI.26}
\end{equation}
The unitary map (TI.11) satisfies
$U(D-1/2)F=-(UF)'$ by differentiation. Integration by parts in the
Fourier transform therefore gives
$\F_+(UH)(t)=\ii t\,g(1/2+\ii t)$. It follows, with the original
$1/(2\pi)$ factor, that
\begin{equation}
 M_1''(0)=\frac1{2\pi}\int_{\mathbb R}t^2|g(1/2+\ii t)|^2\dd t
 =\int_0^\infty |H(x)|^2\dd x=2\int_1^\infty H(x)^2\dd x.
 \tag{TI.27}
\end{equation}
For the last equality, $DR=R(1-D)$ and $Rf_0=f_0$ give
$RH=-H$, that is $H(1/x)=-xH(x)$. Its square gives precisely the
same Jacobian cancellation as before, with the new minus sign retained
before squaring.

Let $H_J$ restrict (TI.26) to $n\le J$ and put
$T_J=2\int_1^\infty H_J(x)^2\dd x$. Then
\begin{equation}
 T_J=\sum_{m,n=1}^J\sum_{r,s=1}^3
 d_rd_s(\pi m^2)^r(\pi n^2)^s
 \frac{\Gamma(r+s+1/2,\pi(m^2+n^2))}
 {\bigl(\pi(m^2+n^2)\bigr)^{r+s+1/2}}. \tag{TI.28}
\end{equation}
There is no coefficient $4$ in this last expression: (TI.26) already
contains its fully differentiated coefficients, and the endpoint
factor $2$ cancels the $1/2$ in (TI.18).
Set
\begin{equation}
 B_a=30\pi H_2(a,\pi)+60\pi^2H_4(a,\pi)+16\pi^3H_6(a,\pi).
 \tag{TI.29}
\end{equation}
The proof of (TI.23), now with degree six in $x$, gives
\begin{equation}
 |M_1''(0)-T_J|
 \le \mathcal T_J:=4B_1B_{J+1}I_{12}(\pi(1+(J+1)^2)). \tag{TI.30}
\end{equation}
In detail, $|H|,|H_J|\le B_1x^6e^{-\pi x^2}$ and
$|H-H_J|\le B_{J+1}x^6e^{-\pi(J+1)^2x^2}$; apply
$|H^2-H_J^2|\le|H-H_J|(|H|+|H_J|)$ and the factor $2$ of (TI.27).

For tensor degree $k\ge1$, let $m_k=w_1^{*k}$ with its total mass
$\mu_1^k$, where $\mu_1=M_1(0)$. Absolute integrability and Fubini give
$Z_k(\theta)=\int e^{\theta u}m_k(u)\dd u=M_1(\theta)^k$.
The function $w_1$ is even by (TI.6), so $Z_k'(0)=M_1'(0)=0$.
The degree-zero monic source polynomial is $1$, of squared norm
$h_0=\mu_1^k$, and its monic degree-one orthogonal polynomial in
the original sum variable $u$ is $u$, of squared norm
$h_1=Z_k''(0)=k\mu_1^{k-1}M_1''(0)$. Therefore the first recurrence
coefficient is exactly
\begin{equation}
 a_1^{(k)}=\frac{h_1}{h_0}
 =k\frac{M_1''(0)}{M_1(0)}. \tag{TI.31}
\end{equation}
The quotient of norms in (TI.31) is the coefficient in the monic
three-term recurrence. It leaves both norms and the measure mass
specified, and performs no replacement of $m_k$ by a probability law.

\section{Certified numerical enclosures and their executable scope}

The accompanying script
\texttt{toda\_theta\_input\_tail\_check\_20260912.py} evaluates
(TI.19), (TI.28), (TI.23), and (TI.30) at $J=8$ with
\texttt{python-flint 0.9.0}, using Arb real balls and Acb rectangular
complex balls at 256-bit precision and one arithmetic thread.
All integer and rational input, including $\theta=1/10$, is supplied
exactly; $\pi$, exponentials, powers, and incomplete gamma functions
are enclosed by those ball operations. The incomplete gamma call has
argument $C$ and parameter $\alpha$, that is
\texttt{C.gamma\_upper(alpha)}. Its branch is the one in (TI.18).
The documented ball and endpoint semantics used here are in
\href{https://python-flint.readthedocs.io/en/latest/arb.html}{the Arb interface}
and
\href{https://python-flint.readthedocs.io/en/latest/acb.html}{the Acb interface}.

The script adds a symmetric error ball whose radius is the upward
rounded upper endpoint of the proved tail bound. Thus the result
encloses the infinite sum, rather than only its finite arithmetic.
The finite real-tilt sum is real by the proved conjugation identity;
the script additionally checks that its imaginary ball contains zero.
The positive real mass enclosure is then used in interval division
for (TI.31). The following are rational decimal intervals:
\begin{align*}
 M_1(0)&\in[
 1.2790072478464851404795335922671932744916,\\[-2pt]
 &\hspace{39mm}1.2790072478464851404795335922671932744919],\\
 M_1(1/10)&\in[
 1.3459071784582497969682205923576298051604,\\[-2pt]
 &\hspace{39mm}1.3459071784582497969682205923576298051607],\\
 M_1''(0)&\in[
 13.0555493025705584353926846581231936824343,\\[-2pt]
 &\hspace{39mm}13.0555493025705584353926846581231936824346],\\
 a_1^{(k)}/k&\in[
 10.2075647534857216888657612522375551749964,\\[-2pt]
 &\hspace{39mm}10.2075647534857216888657612522375551749967].
\end{align*}
Every width is $3\cdot10^{-40}$. The computed positive tail bounds
satisfy, respectively,
\begin{equation}
 {\cal E}_8(0,0)<7\cdot10^{-107},\qquad
 {\cal E}_8(1/10,0)<3\cdot10^{-106},\qquad
 {\cal T}_8<4\cdot10^{-103}. \tag{TI.32}
\end{equation}
The receipt preserves both the displayed decimal intervals and exact
rational dyadic lower and upper endpoints of each resulting ball.
Its decimal endpoints are generated by outward rounding at 40 decimal
places, with one additional unit on each side, and are subsequently
checked by rigorous ball comparisons. Normal and optimized execution
produce identical mathematical records. The requested negative-control
mode inserts the false inequality $a_1^{(k)}/k<1$ and must raise an
error; it uses an explicit checked exception rather than a Python
assertion, so optimization cannot remove that check.

The numerical certificate uses the standard correctness contract of
Arb/Acb ball arithmetic together with the written infinite-tail proof.
It is not inferred from agreement of two numerical precisions. Its
output concerns the three specified infinite analytic-seed quantities
and their recurrence ratio. For $h=1$ the cyclic arithmetic quotient
is $\mathbb C[S]/(1)=0$: the actual quotient map sends each polynomial
to zero, while the positive seed source norm and (TI.31) remain.
For a selected packet $h$, the precise passage from the present input
to $w_h$ is (TI.16), obtained from the globally defined Mellin division
$g/h$. No packet determinant or bound as $k$ tends to infinity is
certified by these four scalar intervals.
\end{document}
